\documentclass[10pt,a4paper]{amsart}
\usepackage[margin=19mm]{geometry}
\usepackage{amsmath,amssymb,mathtools}
\usepackage[scr=rsfs,cal=euler]{mathalfa}
\usepackage{xfrac}
\usepackage[unicode,hidelinks,bookmarks=false]{hyperref}
\newcommand{\ie}{i.e.\ }
\newcommand{\cf}{cf.\ }
\newcommand{\resp}{respectively\ }
\newcommand{\Subsection}{\subsection}
\providecommand{\nobreakdash}{\nobreak\hbox{-}}
\setlength{\emergencystretch}{2em}
\multlinegap0pt
\usepackage{bm}
\usepackage{bbm}
\usepackage{dsfont}
\usepackage{xspace}
\usepackage{cancel}
\usepackage{mathtools}
\usepackage{amsthm}
\makeatletter
\@ifundefined{theorem}{\newtheorem{theorem}{Theorem}}{}
\@ifundefined{proposition}{\newtheorem{proposition}[theorem]{Proposition}}{}
\makeatother
\DeclareMathOperator{\cl}{cl}
\newcommand\Bcal{{\mathcal B}}
\newcommand{\swap}[2]{\xleftrightarrow{#1 \quad #2}}
\title[White's Conjecture for Paving Matroids]{White's Conjecture for Paving Matroids}
\author{Yu-Chuan Yu, Chi Ho Yuen}
\date{Source: arXiv:2510.04163v1 (CC BY 4.0)}
\begin{document}
\maketitle

\noindent\emph{Source and licence.} White's Conjecture for Paving Matroids. Version arXiv:2510.04163v1 (CC BY 4.0), distributed under the Creative Commons Attribution 4.0 International licence, as recorded in the arXiv licence metadata for this exact version and shown on the abstract page https://arxiv.org/abs/2510.04163v1. Full rights notice: licenses/arxiv-2510.04163-CC-BY-4.0.txt.\par\smallskip

\noindent\textit{Editorial scope.} Counted unit: the Proposition whose source label is \texttt{prop:XY\_XY} in the article (source0.tex lines 234--245, source label \texttt{prop:XY\_XY}), delivered together with the article's own complete original proof of it (source0.tex lines 247--309, one \texttt{proof} environment of 63 lines). No mathematics is written, repaired or re-derived: statement and proof are the article's own text line for line. Required context retained uncounted: prop:minor-closed (lines 129--131); prop:SH\_intersection (lines 153--155). Every definition, display and label of those lines is retained unchanged. Omitted from this package and not redistributed: 1--128, 132--152, 156--233, 246--246, 310--838. Nothing inside a retained excerpt is omitted or altered: every retained body is delivered exactly as the article's lines are written, so no omission receipt of any kind accompanies this package. Citations: the retained excerpts carry no citation command at all, so no bibliography entry of the article is transcribed and no reference is redistributed. Reference adaptation: none was needed and none was made; every reference inside the delivered unit resolves inside it, so no unresolved reference or citation is produced. Printed slips kept literal: no printed word, symbol, hypothesis or number of the counted statement or of its proof was altered. Layout and class: the article's own class is \texttt{amsart} with packages including amsaddr, hyperref, amsfonts, amsmath, amssymb, amsthm, amscd, mathtools, mathrsfs, mathabx, graphicx, tikz; the wrapper is the lane's pinned \texttt{amsart} editorial wrapper at 10pt on A4, which restates the theorem-like environments the retained text needs and adds the article's own macro definitions that the retained text needs (\texttt{\textbackslash Bcal}, \texttt{\textbackslash cl}, \texttt{\textbackslash swap}), with extra wrapper packages (bm, bbm, dsfont, xspace, cancel, mathtools). The delivered numbering is the unit's own. Assets: no retained span contains a figure environment, a graphics-inclusion command, a PDF-page inclusion, a table or a TikZ picture; no image file is copied into this package, the package declares no asset dependency and no image file is redistributed with it. Licence: version arXiv:2510.04163v1 of this article is distributed under the Creative Commons Attribution 4.0 International licence, as recorded in the arXiv licence metadata for this exact version and shown on the abstract page https://arxiv.org/abs/2510.04163v1. A full rights notice is delivered at licenses/arxiv-2510.04163-CC-BY-4.0.txt. Characters: every retained excerpt is pure ASCII, so the delivered engine, XeLaTeX with its default Latin Modern text font, reports no missing character.
\begin{proposition} \label{prop:minor-closed}
    The class of paving matroids is minor-closed.
\end{proposition}

\begin{proposition} \label{prop:SH_intersection}
    Let $H,H'$ be two distinct hyperplanes in a paving matroid $M$. Then $|H\cap H'|\leq r-2$.
\end{proposition}

\begin{proposition} \label{prop:XY_XY}
    Let $M$ be a paving matroid and fix a hyperplane $H$ of size $\geq r$. Let $\widetilde{M}$ be the relaxation of $M$ at $H$. Let $X,Y,X',Y'$ be 4 bases of $M$ (hence $\widetilde{M}$) such that there exists an exchange sequence from $(X,Y)$ to $(X',Y')$ in $\widetilde{M}$ as
\begin{equation} \label{eq:deg2}
\begin{array}{c c c c c}
   X   & \swap{a}{s} & X'' & \swap{t}{b} & X'\\
   Y   & \swap{s}{a} & Y'' & \swap{b}{t} & Y'
\end{array}.
\end{equation}


Then there exists an exchange sequence from $(X,Y)$ to $(X',Y')$ in $M$.
\end{proposition} 

\begin{proof}
    We are done if $X'',Y''$ are bases of $M$. On the other hand, they cannot be both of type 0, or else $X,Y\subset X''\cup Y''$ are also of type 0. Therefore we may assume $X''$ is of type 0 and $Y''$ is not. We may further assume $|\{a,b,s,t\}|=4$, otherwise, either $(X,Y)=(X',Y')$, or they only differ by a single symmetric exchange.
    
    By induction on $|E|$ (none of the argument in this paper considers matroids with ground sets larger than $E$), we may assume White's conjecture is true for any smaller paving matroids, thus without loss of generality $E=X\sqcup Y$: we may delete all elements outside of $X\cup Y$ and contract all elements in $X\cap Y$ if it is not, the minor is still paving by Proposition~\ref{prop:minor-closed}, while standard facts on matroid bases with respect to minor operations imply that we can lift an exchange sequence in the minor (guaranteed by White's conjecture) back to the original matroid.

    We write $X=\widetilde{X}+at, Y=\widetilde{Y}+bs$ (hence $X'=\widetilde{X}+bs,Y'=\widetilde{Y}+at$), and call $H=\cl(X'')=\cl(\tilde{X}+st)$ as $H_{X,st}$; note that $a,b\not\in H_{X,st}$, or else $X,Y$ are of type 0 and cannot be bases of $M$.\\
    
    {\bf Step 1}: We may assume $t\not\in C(Y'=\widetilde{Y}+at,s)$, otherwise $\widetilde{Y}+as\in\Bcal$ and
    $$
    \begin{array}{c c c c c}
        \widetilde{X}+at   & \swap{a}{b} & \widetilde{X}+bt & \swap{t}{s} & \widetilde{X}+bs\\
        \widetilde{Y}+bs   & \swap{b}{a} & \widetilde{Y}+as & \swap{s}{t} & \widetilde{Y}+at
    \end{array}
    $$
    is a valid exchange sequence in $M$. Here $\widetilde{X}+bt\in\Bcal(H_{X,st}, b)$.

    Therefore, we may assume $C(\widetilde{Y}+at,s)\subset\widetilde{Y}+as$, which must be an equality as the right hand side is already of size $r$. Consider the hyperplane $H_{Y,as}:=\cl(\widetilde{Y}+as)$ this circuit spans, we have $b,t\not\in H_{Y,as}$, or else $Y,Y'$ cannot be bases of $M$.
    
    Similarly, we may assume $s\not\in C(Y=\widetilde{Y}+bs,t)$, otherwise $\widetilde{Y}+bt\in\Bcal$ and
    $$
    \begin{array}{c c c c c}
        \widetilde{X}+at   & \swap{t}{s} & \widetilde{X}+as & \swap{a}{b} & \widetilde{X}+bs\\
        \widetilde{Y}+bs   & \swap{s}{t} & \widetilde{Y}+bt & \swap{b}{a} & \widetilde{Y}+at
    \end{array}
    $$
    is a valid exchange sequence in $M$ with $\widetilde{X}+as\in\Bcal(H_{X,st}, a)$. So $C(\widetilde{Y}+bs,t)=\widetilde{Y}+bt$ spans a hyperplane $H_{Y,bt}$ that does not contain $a,s$.

    Since $H_{Y,as}\cap H_{Y,bt}\supset\widetilde{Y}$, which is already of size $r-2$, by Proposition~\ref{prop:SH_intersection}, it must be an equality.\\

    {\bf Step 2}: We may assume $\widetilde{X}+ab$ is a circuit, or otherwise
    $$
    \begin{array}{c c c c c}
        \widetilde{X}+at   & \swap{t}{b} & \widetilde{X}+ab & \swap{a}{s} & \widetilde{X}+bs\\
        \widetilde{Y}+bs   & \swap{b}{t} & \widetilde{Y}+st & \swap{s}{a} & \widetilde{Y}+at
    \end{array}
    $$
    is a valid exchange sequence in $M$ with $\widetilde{Y}+st\in\Bcal(H_{Y,as}, t)$. Let $H_{X,ab}$ be the hyperplane it spans. We have $s,t\not\in H_{X,ab}$ and $H_{X,st}\cap H_{X,ab}=\widetilde{X}$ similarly.\\

    {\bf Step 3}: We may assume $E=H_{X,st}\cup H_{X,ab}$, or otherwise we may pick some $p\in E\setminus (H_{X,st}\cup H_{X,ab})$, which is necessarily in $\widetilde{Y}$, and
    $$
    \begin{array}{c c c c c c c}
        \widetilde{X}+at   & \swap{a}{p} & \widetilde{X}+tp & \swap{t}{b} & \widetilde{X}+bp & \swap{p}{s} & \widetilde{X}+bs\\
        \widetilde{Y}+bs   & \swap{p}{a} & \widetilde{Y}-p+abs & \swap{b}{t} & \widetilde{Y}-p+ast & \swap{s}{p} & \widetilde{Y}+at
    \end{array}
    $$
    is a valid exchange sequence in $M$. Here $\widetilde{X}+tp\in\Bcal(H_{X,st},p), \widetilde{X}+bp\in\Bcal(H_{X,ab},p)$, and $\widetilde{Y}-p+abs=(\widetilde{Y}-p+as)+b\in\Bcal(H_{Y,as},b), \widetilde{Y}-p+ast=(\widetilde{Y}-p+as)+t\in\Bcal(H_{Y,as},t)$.

    Similarly, we may assume $E=H_{Y,as}\cup H_{Y,bt}$, or otherwise we may pick some $q\in E\setminus (H_{Y,as}\cup H_{Y,bt})\subset\widetilde{X}$, and
    $$
    \begin{array}{c c c c c c c}
        \widetilde{X}+at   & \swap{q}{s} & \widetilde{X}-q+ast & \swap{t}{b} & \widetilde{X}-q+abs & \swap{a}{q} & \widetilde{X}+bs\\
        \widetilde{Y}+bs   & \swap{s}{q} & \widetilde{Y}+bq & \swap{b}{t} & \widetilde{Y}+tq & \swap{q}{a} & \widetilde{Y}+at
    \end{array}
    $$
    is a valid exchange sequence in $M$. Here $\widetilde{Y}+bq, \widetilde{Y}+tq\in\Bcal(H_{Y,bt},q)$, and $\widetilde{X}-q+ast=(\widetilde{X}-q+st)+a\in\Bcal(H_{X,st},a), \widetilde{X}-q+abs=(\widetilde{X}-q+ab)+s\in\Bcal(H_{X,ab},s)$.\\

    {\bf Step 4}: For each pair of distinct hyperplanes chosen from $H_{X,st}, H_{X,ab}, H_{Y,as}, H_{Y,bt}$, we consider the size of their intersection $|H_{\cdot,\cdot}\cap H_{\cdot,\cdot}|$ and sum over the $\binom{4}{2}=6$ terms. It can be seen that
    \begin{enumerate}
        \item each of $a,b,s,t$ appears in exactly 2 hyperplanes and contributes $\binom{2}{2}=1$ to the sum;
        \item each element in $\widetilde{X}\cup\widetilde{Y}$ appears in exactly 3 hyperplanes and contributes $\binom{3}{2}=3$ to the sum: for example, suppose $x\in\widetilde{X}$ thus $x\in H_{X,st}, H_{X,ab}$, by Step 3 we know that $x$ belongs to exactly one of $ H_{Y,as}, H_{Y,bt}$ as well.
    \end{enumerate}
    Therefore the sum is $3(|\widetilde{X}|+|\widetilde{Y}|)+4=3\cdot 2(r-2)+4=6r-8$, but on the other hand, the sum is at most $6(r-2)=6r-12$ by Proposition~\ref{prop:SH_intersection}, a contradiction.
\end{proof}

\end{document}
